\documentclass[10pt,a4paper]{amsart}
\usepackage{amsmath,amssymb,amsthm,graphicx}
\usepackage{verbatim}
\newenvironment{reference}{\comment}{\endcomment}
%\newenvironment{reference}{}{}
\newenvironment{slogan}{\comment}{\endcomment}
\newenvironment{history}{\comment}{\endcomment}

% For commutative diagrams we use Xy-pic
\usepackage[all]{xy}

% We use 2cell for 2-commutative diagrams.
\xyoption{2cell}
% Theorem environments.
%
\theoremstyle{plain}
\newtheorem{theorem}[subsection]{Theorem}
\newtheorem{proposition}[subsection]{Proposition}
\newtheorem{lemma}[subsection]{Lemma}

\theoremstyle{definition}
\newtheorem{definition}[subsection]{Definition}
\newtheorem{example}[subsection]{Example}
\newtheorem{exercise}[subsection]{Exercise}
\newtheorem{situation}[subsection]{Situation}

\theoremstyle{remark}
\newtheorem{remark}[subsection]{Remark}
\newtheorem{remarks}[subsection]{Remarks}

\numberwithin{equation}{subsection}

% Macros
%
\def\lim{\mathop{\mathrm{lim}}\nolimits}
\def\colim{\mathop{\mathrm{colim}}\nolimits}
\def\Spec{\mathop{\mathrm{Spec}}}
\def\Hom{\mathop{\mathrm{Hom}}\nolimits}
\def\Ext{\mathop{\mathrm{Ext}}\nolimits}
\def\SheafHom{\mathop{\mathcal{H}\!\mathit{om}}\nolimits}
\def\SheafExt{\mathop{\mathcal{E}\!\mathit{xt}}\nolimits}
\def\Sch{\mathit{Sch}}
\def\Mor{\mathop{\mathrm{Mor}}\nolimits}
\def\Ob{\mathop{\mathrm{Ob}}\nolimits}
\def\Sh{\mathop{\mathit{Sh}}\nolimits}
\def\NL{\mathop{N\!L}\nolimits}
\def\CH{\mathop{\mathrm{CH}}\nolimits}
\def\proetale{{pro\text{-}\acute{e}tale}}
\def\etale{{\acute{e}tale}}
\def\QCoh{\mathit{QCoh}}
\def\Ker{\mathop{\mathrm{Ker}}}
\def\Im{\mathop{\mathrm{Im}}}
\def\Coker{\mathop{\mathrm{Coker}}}
\def\Coim{\mathop{\mathrm{Coim}}}

% Boxtimes
%
\DeclareMathSymbol{\boxtimes}{\mathbin}{AMSa}{"02}

%
% Macros for moduli stacks/spaces
%
\def\QCohstack{\mathcal{QC}\!\mathit{oh}}
\def\Cohstack{\mathcal{C}\!\mathit{oh}}
\def\Spacesstack{\mathcal{S}\!\mathit{paces}}
\def\Quotfunctor{\mathrm{Quot}}
\def\Hilbfunctor{\mathrm{Hilb}}
\def\Curvesstack{\mathcal{C}\!\mathit{urves}}
\def\Polarizedstack{\mathcal{P}\!\mathit{olarized}}
\def\Complexesstack{\mathcal{C}\!\mathit{omplexes}}
% \Pic is the operator that assigns to X its picard group, usage \Pic(X)
% \Picardstack_{X/B} denotes the Picard stack of X over B
% \Picardfunctor_{X/B} denotes the Picard functor of X over B
\def\Pic{\mathop{\mathrm{Pic}}\nolimits}
\def\Picardstack{\mathcal{P}\!\mathit{ic}}
\def\Picardfunctor{\mathrm{Pic}}
\def\Deformationcategory{\mathcal{D}\!\mathit{ef}}
\usepackage{cleveref}
\crefname{originallemma}{original lemma}{original lemmas}
\Crefname{originallemma}{Original lemma}{Original lemmas}
\makeatletter
\newlabel{lemma-karoubian}{{09SH}{source}{}{sourcepointer-1}{}}
\newlabel{lemma-karoubian@cref}{{[originallemma][0][]09SH}{[1][0][]source}{}{}{}}
\makeatother
\title{Countable products and kernels of split maps make a preadditive category Karoubian}
\author{The Stacks Project Authors}
\date{stacks-project:homology.tex:local-snapshot-20261012; GNU FDL 1.2 or later}
\begin{document}
\maketitle
\noindent\textit{Editorial note.} All delivered mathematical text is the author's own, copied from the pinned Stacks Project chapter file homology.tex; no statement or proof step was rewritten, reordered or removed.\par\medskip
\noindent\textit{Editorial note.} The proof is the author's complete argument for the first assertion; the closing sentence 'The proof of (2) is dual (using the dual definition of a Karoubian category, namely condition (2) of Lemma lemma-karoubian).' is the author's own duality remark on the part written out in full inside the counted span, and it is preserved verbatim.\par\medskip
\noindent\textit{Editorial note.} Every labelled reference occurring inside the counted spans is carried in assembly.external\_references as a non-load-bearing navigation pointer (same-chapter ones at their definition site, cross-chapter ones at the invocation inside the counted span), so no load-bearing mathematics is replaced by a stub. No bibliography citation occurs inside the counted spans.\par\medskip
\noindent\textit{Editorial note.} The same-name check for this card scanned every staged Stacks chapter of this repository -- all sources/stacks-*/pinned-original-source.tex files, the card's own chapter excluded -- and the labels found are recorded in dedup.same\_name\_cross\_chapter on the card.\par\medskip
\noindent\textit{Editorial note.} The preamble is the chapter's own native preamble excerpt already used for this source: the comment-type reference, slogan and history environments (native preamble lines 8-13), the Xy-pic setup (native preamble lines 14-18) and the native theorem declarations with the mathematical macros (native preamble lines 154-217, the actual last line of the file).\par\medskip
\noindent\textit{Allowed standard prerequisites (editorial).}\begin{itemize}
\item Karoubian categories. Same chapter, Lemma lemma-karoubian. Supplies the two equivalent descriptions of a Karoubian category, in particular the dual description invoked for the second assertion of the proof.
\item Kernels as representing objects. Same chapter, the definition of a kernel and its description by a representable functor. Supplies the observation that the functor sending an object W to the kernel of the induced endomorphism of the Hom-set is representable exactly when the idempotent has a kernel.
\end{itemize}\medskip
\ifdefined\hypertarget\hypertarget{sourcepointer-1}{}\fi
\noindent\textit{Source reference.} Original lemma 09SH: sources/stacks-homology/pinned-original-source.tex, source lines 743--743. Navigation pointer invoked at proof line 803; the counted span contains the pointer only and the cited result is not reproduced inside it, so the pointer is registered as a non-load-bearing navigation stub in the established form for Stacks material. No bibliography citation is involved.\par\medskip
\begin{lemma}
\label{lemma-projectors-have-images}
Let $\mathcal{D}$ be a preadditive category.
\begin{enumerate}
\item If $\mathcal{D}$ has countable products and kernels of maps which
have a right inverse, then $\mathcal{D}$ is Karoubian.
\item If $\mathcal{D}$ has countable coproducts and cokernels of
maps which have a left inverse, then $\mathcal{D}$ is Karoubian.
\end{enumerate}
\end{lemma}

\begin{proof}
Let $X$ be an object of $\mathcal{D}$ and let $e : X \to X$ be an idempotent.
The functor
$$
W \longmapsto \Ker(
\Mor_\mathcal{D}(W, X)
\xrightarrow{e}
\Mor_\mathcal{D}(W, X)
)
$$
is representable if and only if $e$ has a kernel. Note that for any
abelian group $A$ and idempotent endomorphism $e : A \to A$ we have
$$
\Ker(e : A \to A)
= \Ker(\Phi :
\prod\nolimits_{n \in \mathbf{N}} A
\to
\prod\nolimits_{n \in \mathbf{N}} A
)
$$
where
$$
\Phi(a_1, a_2, a_3, \ldots) = (ea_1 + (1 - e)a_2, ea_2 + (1 - e)a_3, \ldots)
$$
Moreover, $\Phi$ has the right inverse
$$
\Psi(a_1, a_2, a_3, \ldots) =
(a_1, (1 - e)a_1 + ea_2, (1 - e)a_2 + ea_3, \ldots).
$$
Hence (1) holds. The proof of (2) is dual (using the dual definition
of a Karoubian category, namely condition (2) of
Lemma \ref{lemma-karoubian}).
\end{proof}

\end{document}
